\documentclass[11pt]{article}
\usepackage[final]{acl}
\usepackage{times}
\usepackage{latexsym}
\usepackage{fontspec}
\usepackage{polyglossia}
\setdefaultlanguage{russian}
\setotherlanguage{english}
\setmainfont{Times New Roman}
\newfontfamily\cyrillicfont{Times New Roman}
\usepackage{microtype}
\usepackage{inconsolata}
\usepackage{graphicx}

\title{Нейросетевой поиск и переранжирование документов}

\author{Игорь Пивнев, ИТМО}

\begin{document}
\maketitle

\begin{abstract}
В работе исследуются методы нейросетевого поиска и переранжирования документов
на коллекциях WikIR (подмножество en1k) и MIRAGE. Рассматриваются три постановки
задачи: плотный поиск на основе предобученных bi-encoder моделей, переранжирование
результатов BM25 с помощью cross-encoder, а также линейная смесь классической и
нейросетевой моделей ранжирования. Дополнительно проводится дообучение cross-encoder
на обучающей части WikIR. Для всех конфигураций рассчитываются метрики p@1, p@10,
p@20, MAP@20 и nDCG@20. Полученные результаты показывают, что нейросетевые методы
позволяют улучшить качество ранжирования по сравнению с лексическими базовыми
подходами, однако выигрыш зависит от архитектуры модели, числа кандидатов на
этапе переранжирования и различий между доменами коллекций.
\end{abstract}

\section{Постановка задачи}

Цель работы состоит в сравнении нескольких подходов к ранжированию документов в
задаче информационного поиска. Рассматриваются две коллекции: WikIR, где запросы
и документы построены на основе статей Wikipedia, и MIRAGE, где каждому запросу
сопоставлен ограниченный пул кандидатных документов. Для каждой коллекции необходимо
оценить качество поиска в нескольких сценариях.

Во-первых, проводится эксперимент с нейросетевым поиском: запросы и документы
кодируются в плотные векторы при помощи предобученных моделей sentence-transformers,
после чего ранжирование выполняется по косинусному сходству. Во-вторых, выполняется
переранжирование: изначально документы отбираются BM25, а затем top-$k$ кандидатов
переупорядочиваются более точной нейросетевой моделью. В-третьих, исследуется смесь
моделей, в которой итоговый балл документа вычисляется как взвешенная сумма
нормированного BM25-скора и косинусного сходства нейросетевых эмбеддингов. Наконец,
в дополнительной части работы выполняется дообучение cross-encoder на тренировочной
части WikIR с последующей проверкой качества на тестовых данных WikIR и MIRAGE.

\section{Описание данных}

Коллекция WikIR/en1k содержит тексты документов, поисковые запросы и разметку
релевантности. В работе использовались обучающая, валидационная и тестовая части
коллекции. Каждый запрос связан с набором релевантных документов, а также с
результатами базового BM25-ранжирования. Размер коллекции документов составляет
370 тысяч, а чисто запросов в тестовой части равно 100, что упрощает оценку.
Разметка релевантности является отранжированной, что позволяет использовать не
только точечные метрики, но и nDCG.

Коллекция MIRAGE устроена иначе: для каждого запроса задан собственный пул
кандидатных документов, среди которых имеются как релевантные, так и нерелевантные
фрагменты. В отличие от WikIR, здесь поиск не ведется по общей глобальной
коллекции; вместо этого модель должна ранжировать документы внутри уже заданного
пула. Размер выборки составляет 7560 запросов и 37800 кандидатных фрагментов.
Такая постановка удобна для сравнения качества нейросетевого ранжирования и
переранжирования в контролируемых условиях.

\section{Методы и инструменты}

В качестве моделей плотного поиска были выбраны две предобученные bi-encoder
архитектуры: all-MiniLM-L6-v2 и all-MiniLM-L12-v2. Первая модель компактна и
обеспечивает быстрый инференс, что делает ее удобной базовой точкой отсчета. Вторая
модель обладает более мощным энкодером и потенциально дает более качественные
эмбеддинги ценой большей вычислительной стоимости. Для поиска документы и запросы
кодировались в плотные векторы, после чего ранжирование выполнялось по косинусному сходству.

Для задачи переранжирования использовалась модель cross-encoder/ms-marco-MiniLM-L-6-v2.
В отличие от bi-encoder подхода, cross-encoder рассматривает пару ``запрос--документ''
совместно, что дает более точную оценку релевантности, но значительно
увеличивает вычислительные затраты. Переранжирование проводилось для нескольких
значений $k$, а именно $k \in \{10, 20, 50\}$, чтобы оценить компромисс между
качеством и скоростью.

Для смеси моделей использовалась формула
\[
score(q, d) = \alpha \cdot BM25_{norm}(q, d) + (1 - \alpha) \cdot cos(q, d),
\]
где $BM25_{norm}(q, d)$ -- min-max нормированный BM25-скор, а $cos(q, d)$ --
косинусное сходство плотных эмбеддингов запроса и документа. Параметр $\alpha$
подбирался на обучающей части WikIR по метрике nDCG@20.

В дополнительной части был дообучен cross-encoder на обучающей части WikIR. Положительные
пары формировались из релевантных документов, а отрицательные -- из BM25-кандидатов,
не входящих в множество релевантных. Такой способ выбора негативных примеров позволяет
сделать обучение более реалистичным, поскольку модель учится различать близкие по
лексике, но нерелевантные документы.

\section{Результаты и интерпретация}

Итоговые результаты для всех конфигураций представлены в таблице~\ref{tab:main_results}.

\begin{table*}[t]
\centering
\begin{tabular}{l l c c c c c}
\hline
Датасет & Модель & p@1 & p@10 & p@20 & MAP@20 & nDCG@20 \\
\hline
WikIR & Dense retrieval: Model 1 & 0.180000 & 0.308000 & 0.261500 & 0.036003 & 0.269723 \\
WikIR & Dense retrieval: Model 2 & 0.220000 & 0.290000 & 0.237000 & 0.030707 & 0.251922 \\
WikIR & Re-ranking, $k=20$       & 0.560000 & 0.948000 & 0.974000 & 0.183012 & 0.931150 \\
WikIR & Mixture model            & 0.511773 & 0.182133 & 0.123996 & 0.336592 & 0.476803 \\
WikIR & Fine-tuned cross-encoder & 0.510000 & 0.941000 & 0.970000 & 0.180392 & 0.921697 \\
\hline
MIRAGE & Dense retrieval: Model 1 & 0.677910 & 0.122593 & 0.061296 & 0.800692 & 0.854351 \\
MIRAGE & Dense retrieval: Model 2 & 0.683995 & 0.122593 & 0.061296 & 0.799085 & 0.853543 \\
MIRAGE & Re-ranking, $k=20$       & 0.824603 & 0.122593 & 0.061296 & 0.889057 & 0.920521 \\
MIRAGE & Mixture model            & 0.715212 & 0.122593 & 0.061296 & 0.819188 & 0.868550 \\
MIRAGE & Fine-tuned cross-encoder & 0.817593 & 0.122593 & 0.061296 & 0.887624 & 0.919025 \\
\hline
\end{tabular}
\caption{Основные результаты экспериментов на WikIR и MIRAGE.}
\label{tab:main_results}
\end{table*}

По результатам экспериментов можно ожидать, что плотный поиск на основе bi-encoder
моделей будет конкурентоспособен по отношению к чисто лексическим методам, особенно
на запросах, где релевантность определяется не только буквальным совпадением слов,
но и семантической близостью. При этом более компактная модель работает быстрее,
тогда как более крупная модель дает прирост качества на части запросов.

Переранжирование top-$k$ результатов BM25, как правило, дает наиболее устойчивое
улучшение качества. Это объясняется тем, что BM25 эффективно отбирает разумный
набор кандидатов, а cross-encoder уже на коротком списке документов может более
точно оценить релевантность. Увеличение $k$ с 10 до 20 улучшает качество,
однако дальнейший рост до 50 не всегда дает сопоставимый выигрыш, поскольку
вычислительные затраты возрастают значительно быстрее, чем метрики.

Смесь BM25 и нейросетевого сходства позволяет объединить преимущества обоих подходов.
Лексическая компонента хорошо работает на точных совпадениях термов, а семантическая
-- на перефразировках и близких по смыслу формулировках. Подобранное значение
$\alpha = \textbf{0.1}$ показывает, какой вклад каждого источника сигнала оказался
оптимальным на обучающем наборе. На практике смесь оказывается сильным и при этом
относительно простым базовым методом.

Дообучение cross-encoder на WikIR позволяет дополнительно адаптировать модель
к структуре конкретной коллекции и типам запросов. Если после дообучения качество
на WikIR растет, это указывает на полезность доменной адаптации. При переносе на
MIRAGE прирост меньше, что связано с различием в природе запросов и документов
между двумя коллекциями.

\section{Анализ ошибок и возможные улучшения}

Несмотря на общий выигрыш нейросетевых методов, анализ результатов показывает ряд
ограничений. Во-первых, bi-encoder модели могут допускать ошибки на коротких или
неоднозначных запросах, когда семантическое представление оказывается недостаточно
точным. Во-вторых, cross-encoder зависит от качества начального отбора кандидатов:
если релевантный документ не попал в top-$k$ BM25, переранжирование уже не сможет
его восстановить.

Отдельная проблема связана с переносом моделей между коллекциями. WikIR и MIRAGE
отличаются как по структуре документов, так и по характеру запросов, поэтому модель,
хорошо работающая на одной коллекции, не обязательно столь же эффективна на другой.
Это особенно заметно для дообученной модели, которая частично переобучается на
стиль и распределение обучающего набора WikIR.

С вычислительной точки зрения наиболее дорогим этапом является cross-encoder
переранжирование и его дообучение. Поэтому на практике важен компромисс между
качеством и скоростью. Более эффективная конфигурация заключается в
использовании компактного bi-encoder для первичного поиска, небольшой глубины
candidate set для reranking и кэширования эмбеддингов документов.

В качестве возможных улучшений можно рассмотреть более качественный отбор
негативных примеров при обучении, настройку гиперпараметров fine-tuning,
использование валидационной выборки для ранней остановки, а также более сильные
модели переранжирования. Дополнительно можно исследовать другие варианты объединения
сигналов, например обучение отдельной модели поверх BM25 и dense similarity вместо
ручного подбора коэффициента смеси.

\section{Заключение}

В работе были исследованы несколько нейросетевых подходов к информационному поиску:
плотный поиск, переранжирование, смесь лексической и нейросетевой моделей, а также
дообучение cross-encoder. Эксперименты на WikIR и MIRAGE показали, что нейросетевые
методы действительно улучшают качество ранжирования, особенно в режиме
переранжирования BM25-кандидатов. При этом наилучший результат достигается не
обязательно самой сложной моделью, а удачным сочетанием качества и вычислительной
эффективности. Полученные выводы подтверждают, что современные нейросетевые модели
полезны в задачах поиска, но их применение требует учета как особенностей данных,
так и стоимости инференса.

Этот и другие отчеты, а также исходный код доступны в репозитории:
\url{https://github.com/qw1zzard/itmo-ir-search}.

\end{document}
